% Editar este fichero y rellenar los datos

% título y subtítulo del TFM en español
\newcommand{\tituloTFM}{
    Generación automática de informes clínicos a partir de dictado médico
}
\newcommand{\subtituloTFM}{
    % usando PLN y ASR para extraer entidades clínicas, generación de informes clínicos 
    % estructurados con LLMs y estándar FHIR
}
% palabras clave en español
\newcommand{\palabrasClave}{PLN, ASR, LLMs, FHIR, Generación de informes clínicos, SNOMED CT}

% título y subtítulo del TFM en inglés
\newcommand{\titleMT}{Automatic generation of clinical reports from medical dictation}
\newcommand{\subtitleMT}{
    % Extraction of clinical entities, generation of clinical reports using LLMs with structured output (FHIR)
}
% palabras clave en inglés
\newcommand{\keywords}{ASR, PLN, LLMs, Clinical report generation, SNOMED CT, FHIR, HCE}

% Datos del estudiante (nombre y DNI o equivalemente)
\newcommand{\estudiante}{Eduardo Enrique Suárez Castro}
\newcommand{\dni}{Z3632769Q}

% Datos del tutor (nombre, área y departamento)
\newcommand{\tutorA}{Francisco Javier Melero Rus}
\newcommand{\areaA}{Lenguajes y Sistemas Informáticos}
\newcommand{\departamentoA}{Lenguajes y Sistemas Informáticos}

% Datos del co-tutor (nombre, área y departamento), si no hay co-tutoror declarar vacío
\newcommand{\tutorB}{}

 % poner el nombre de la imagen como logotipo o declarar vacío para ningún logotipo:
% \newcommand{\logoTFM}{logo}
\newcommand{\logoTFM}{}

% versión de la memoria
\newcommand{\version}{1.0}


